\section{Discussion, Limitations, and Future Work}\label{sec:discussion}

\subsection{Regime Boundaries as Intrinsic System Transitions}
A regime boundary, in our usage, is the operating point at which one term of the latency decomposition begins to dominate the others.
In the capacity-limited regime, demand exceeds fast-tier capacity and locality cannot help regardless of the coordination strategy.
In the transfer-limited regime, sustained storage bandwidth is the binding constraint: once cross-tier traffic saturates it, orchestration adds latency rather than offsetting it.
Between the two lies the coordination-dominated regime, in which compute, locality, transfer, and synchronization are of comparable magnitude; this is the only regime in which orchestration materially reshapes behavior, and outside it the coordination overhead is not recovered.

\subsection{Fundamental Limits and Inconsistent Prior Results}
This regime view also clarifies apparent disagreements in the literature.
Hierarchical orchestration is not a universal gain, and no implementation makes it one.
Sustained bandwidth is a hard ceiling that coordination cannot exceed, and highly irregular access patterns can remove the benefit on their own: once prediction and coordination cost more than selective transfer saves, the net effect becomes negative. Both failure modes are intrinsic properties of multi-tier memory rather than implementation defects.
The same lens explains why offloading, swapping, and cache-centric studies report inconsistent results---large gains in some configurations and zero or negative returns in others. Viewed through the regime model, those techniques implicitly assume the coordination-dominated regime; deployed outside it, physical constraints dominate and previously amortized overheads become visible.

\subsection{RAMSES as an Experimental Embodiment}
RAMSES is, in this sense, the regime view made operational.
It treats VRAM, DRAM, and NVMe as one hierarchy and ties the gains we measured back to the four-term latency decomposition (compute, locality, transfer, synchronization).
For deployment, RAMSES drops into PyTorch-based stacks via \texttt{LD\_PRELOAD} and configuration files---no application changes required---and can also be packaged as a lightweight container or sidecar~\cite{swapAdvisor2020,kwon2023pagedattention}.

\subsection{What the Evidence Contributes}
Preallocation, NUMA placement, asynchronous I/O, and eviction are not claimed
as new in isolation. The research contribution is the explicit connection
between observable dimensionless pressures, a reproducible regime-to-action
mapping, and a counterfactual evaluation that seeks both positive and negative
operating regions. A faster implementation alone would establish an engineering
point result; agreement among the predicted boundary, policy-off ablation, and
measured response instead supports an explanation of when coordination is
worthwhile.

That explanation is the transferable artifact. A new device, model, or storage
tier need not inherit our numerical thresholds: it can measure its own capacity
and transfer terms, locate its operating point, and test the same boundary
prediction. Conversely, failure to reproduce the predicted ordering, or a
benefit in a supposedly capacity- or bandwidth-limited region, would falsify
the regime account rather than being dismissed as a tuning problem.

\subsection{Generality and Energy Implications}
The experiments here are on LLM and vision-model inference, but regime-dependent behavior is a general property of multi-tier memory.
Anywhere working sets exceed fast-tier capacity and execution depends on cross-tier movement plus synchronization, the same regime transitions appear.
On constrained industrial edge boxes---where $C_{\mathrm{VRAM}} \ll C_{\mathrm{DRAM}}$ (the GPU fast-tier is the binding constraint) or where secondary storage bandwidth is tight---you sit intrinsically near the capacity boundary ($\alpha \approx 1$) or the transfer boundary ($\beta \approx 1$), and instabilities amplify. That matches our observations.
On energy: we measured synchronized whole-node energy (GPU via NVML plus CPU package and DRAM via RAPL).
In real industrial deployments, rack-level PDU or three-phase meters additionally pick up NVMe-rail, cooling, and PSU loads, so certified absolute numbers will be higher. The relative energy--latency trade-offs that come from regime transitions, though, hold up~\cite{huang2019gpipe,stojkovic2025dynamollm}.

\subsection{Limitations and Future Work}
\label{sec:future_work}

We note several limitations explicitly.
First, the evaluation is single-node. Multi-node and cross-fabric deployments---NVLink or InfiniBand clusters, for example---may introduce extra communication-induced regime shifts that we have not characterized here.
Second, the 72-hour trace is synthetic and the available artifact lacks factory reference samples. We therefore withdraw distribution-fidelity claims; real traces may require a different arrival model and calibration procedure.
Third, energy measurement integrates GPU (NVML) with CPU-package and DRAM (RAPL) on a shared clock, cross-validated against an external PDU (Kendall $\tau=1.0$). NVMe-rail power and facility-level overheads (cooling, PSU inefficiencies) are not captured, so absolute certified accounting remains future work.

Three directions look natural from where we ended up.
Extending the regime-aware model to multi-node infrastructures and to heterogeneous accelerators (Neural Processing Units, NPUs; Field-Programmable Gate Arrays, FPGAs).
Building publicly reproducible industrial trace benchmarks so other groups can rerun this kind of analysis without proprietary data.
Adding calibrated power instrumentation for absolute industrial energy certification under IEC~62053.
Of these, the most interesting open question is probably adaptive online policy: an agent that infers the operating regime in real time and tunes orchestration as the workload moves, rather than relying on a static configuration~\cite{stojkovic2025dynamollm}.
